\section{Conclusion}
\label{sec:conclusion}

We set out to measure whether Video-LLMs can explain dashcam collisions faithfully, and we separated the two ways they can fail.
Zero-shot models omit: on CCD they leave out nearly half of the forensic content and often never mention the crash.
Domain adaptation with QLoRA removes the omissions almost entirely and lifts every lexical metric by a wide, statistically unambiguous margin.
It does not touch hallucination.
The adapted model asserts every field on every video, one in five of those assertions contradicts the ground truth, one in five explanations names a vehicle of the wrong colour, and every confirmed no-crash video receives a fabricated collision under the decoding variants that score best on lexical metrics.

We also showed that the adapted model's crash-window predictions are a reproduction of the training-set marginal---two windows, no correlation with the true onset, beaten by a constant---and that this arises because the model is never told what time a frame is.
Temporal IoU reported without a prior-only baseline would have concealed this, and we suggest that such baselines become standard in accident-VLM evaluation.

Finally, we tested the natural training-free remedy and found that it does not work here, and we explained why: a reversed-frame negative preserves every order-invariant cue, so it cannot reach the errors that dominate---wrong colours, wrong labels, and a prior that does not depend on order to begin with.
Mild contrast is harmless and nudges causal phrasing in the right direction; strong contrast makes both omission and hallucination significantly worse.

The constructive message is about what to fix next.
Omission is a data and format problem, and it is solved.
Hallucination on this task has three identifiable sources---no time signal, a prompt that presupposes a crash, and a supervised objective that rewards confident assertion---and each suggests a specific intervention: time-stamped inputs or time tokens, a no-crash-aware prompt and target, and a training signal that penalises unsupported claims rather than rewarding their presence.
We release all outputs, per-video scores, and evaluation code so that those interventions can be measured against the same bar.
